\begin{frame}{Du modèle relationnel au NoSQL}

    Remplacer diagramme de classes par diagramme UML de bases de données

    \begin{figure}[htb]
        \resizebox{\textwidth}{!}{
        \begin{tikzpicture}
            \tikzstyle{descript}=[text=black,align=center,minimum height=1.8cm,align=center,outer sep=0pt]
            \tikzstyle{activity}=[align=center,outer sep=1pt]

            %% Coordinates
            \coordinate (O) at (-1,0); % Origin
            \coordinate (F) at (15,0); % End
            \coordinate (P1) at (0,0);
            \coordinate (P2) at (2,0);
            \coordinate (P3) at (12,0);
            \coordinate (P4) at (14.9,0);

            %% Filled regions
            \fill[color=violet!70]
                (P2) rectangle (P4 |- 1, 1); % RDBMS
            \fill[color=vert]
                (P3 |- 1, 1) rectangle (14.9, 2); % NoSQL (starting from 2010)
            % \fill[color=vert,pattern={mylines[size=5pt,line width=.8pt,angle=45]}]
            %     (P2 |- 1, 1) rectangle (14.9, 2); % NoSQL (starting from 2010)

            \draw [decorate,decoration={brace,amplitude=6pt},text width=2.5cm,align=center,color=violet]($(P2)+(0,3)$) -- ($(P4)+(0,3)$) node [black,midway,above=6pt] {\textcolor{violet}{RDBMS}};
            \draw [decorate,decoration={brace,amplitude=6pt},text width=2.5cm,align=center,color=vert]($(P3)+(0,2)$) -- ($(P4)+(0,2)$) node [black,midway,above=6pt] {\textcolor{vert}{NoSQL}};

            %% Arrow
            \draw[->] (O) -- (F);

            %% Ticks
            \foreach \x in {0,2,...,14}
                \draw(\x cm,3pt) -- (\x cm,-3pt);

            %% Labels
            \foreach \i \j in {0/1950,2/1960,4/1970,6/1980,8/1990,10/2000,12/2010,14/2020}{
                \draw (\i,0) node[below=3pt] {\j} ;
            }
        \end{tikzpicture}
        }
    \end{figure}
\end{frame}
